\section{Conclusion}
\label{sec:conclusion}

We benchmark VLMs on two robot-supervision capabilities and arrive at three findings.
First, zero-shot state reasoning is weak but non-trivial: accuracy stays only marginally
above random across model families, scale helps within a family, but even 30B+ models
plateau well short of reliable supervision. Second, chain-of-thought prompting hurts
rather than helps: two independent models drop below random on \texttt{task\_success}
with CoT, suggesting free-form reasoning chains introduce hallucinated conclusions on
perceptual tasks rather than improving them. Third, fine-tuning gains are real but mostly
coarse: LoRA lifts Qwen3-4B accuracy by +10.9pp, but roughly 80\% of that gain decomposes
into a generalizable out-of-distribution scene detector and a label-position bias that is
in fact a regression for the most critical failure-detection cases. Taken together, these
results suggest current VLMs are not yet ready for reliable robot-state supervision out
of the box, and that fine-tuning gains must be decomposed by question type and answer
class before they can be trusted for deployment.

\section{Limitations \& Future Work}
\label{sec:limitations}

\textbf{Limitations.} Manual annotation limits scale: 47 scenes is enough for a
controlled study but too small to draw strong statistical conclusions about
model-family differences; an automated annotation pipeline would let this scale up.
Train and test splits are drawn from the same Robo2VLM-1 distribution, so our results do
not test sim-to-real or cross-dataset generalization. Some OXE episodes cannot be traced
back to their raw source, limiting reproducibility for those questions. Two runs were
incomplete at the time of writing: the Qwen3-4B CoT run (814 of 6224 questions) and the
Qwen3-8B multi-view run (899 of 1367 source groups, $\sim$66\%) -- the corresponding
numbers may shift once these complete. Finally, LoRA fine-tuning was only run on
Qwen3-4B, so it is unknown whether the CnBD-generalization and IPI-collapse pattern holds
at larger scale.

\textbf{Future work.} We plan to use the Robo2VLM-1 pipeline to generate a larger-scale
dataset for more robust cross-difficulty testing; to test a list-position hypothesis for
\texttt{phase\_success} by swapping \texttt{label\_phase} text for the same-position
phase from a different scene (if the answer is unchanged, this would confirm a
position-only cue rather than temporal grounding); to run LoRA on a larger model
(Qwen3-8B) to check whether the gain-decomposition pattern holds at scale; and to run an
error taxonomy on the CoT transcripts to confirm the ``correct description, wrong
conclusion'' failure mode with concrete examples.
